\paragraph{Strengthening to mean absolute error.}
In the rigidity theorem one may replace $|d_n|^2$ by $|d_n|$.
Here is a proof using the idempotent theorem in harmonic analysis.
Suppose the mean absolute error is bounded by $C$ on windows of
length $H\to\infty$. Each $|d_n|$ on such a window is at most $CH$,
so the preceding Vandermonde argument, applied near its end, still
eliminates unmatched outer and inner points. Let $\Gamma$ be the
countable subgroup of the circle generated by the remaining points $w$,
given the discrete topology. Its dual $G$ is compact and metrizable.
Write $g_n\in G$ for the character $w\mapsto w^n$.

Take triangular nonnegative weights $t_n$ of total mass one on the
window with its last two positions omitted. Explicitly, normalize the
sequence $\min(i+1,H-2-i)$ for $0\le i\le H-3$, with zero values
outside. They satisfy $\sup t_n=O(1/H)$ and
$\sup|t_{n+1}-t_n|=O(1/H^2)$. The measures
$\sum_n t_nd_n\delta_{g_n}$ have bounded variation, so converge
weak-star on a subsequence to a finite measure $\nu$.
For each $z\in\Gamma$, Abel summation shows that the $t_n$-weighted
average of $(w/z)^n$ tends to zero for $w\ne z$ (the total variation
of the weights is $O(1/H)$). The weighted average of $a_nz^{-n}$
therefore tends to $c_z$, taking $c_z=0$ for points outside the
spectrum. Shifting indices as above introduces an error at most
$O(H^{-2}\sum_{n\in I}|d_n|)=O(C/H)$. Hence
$\widehat\nu(z)=-m_z$ for $z\ne1$. Subtracting
$\widehat\nu(1)$ times Haar probability sets the trivial coefficient
to zero.

The resulting Fourier transform is integer-valued and bounded.
Choose a polynomial taking value zero at zero and one at every
nonzero integer in its finite range. Applying its convolution
polynomial to the measure gives a finite measure with Fourier
transform $1_D$, where $D=\{w:m_w\ne0\}$. The
Cohen--Helson--Rudin idempotent theorem puts $D$ in the coset ring
of $\Gamma$. See B.~Green and T.~Sanders,
\href{https://doi.org/10.4007/annals.2008.168.1025}{\emph{A quantitative version of the
idempotent theorem in harmonic analysis}}.

An infinite coset-ring subset of a subgroup of the circle has closure
the entire circle in its usual topology. To verify this, express a
Boolean combination of finitely many cosets as a union of atoms.
Intersecting the positive cosets in one atom yields a single coset
of a subgroup $T$ (possibly the whole group). Split it into finitely
many smaller cosets to settle all the excluded cosets whose subgroups
intersect $T$ in finite index, using the intersection of those
finite-index subgroups. Each surviving piece is a coset of a subgroup
$T_0$ with finitely many cosets of infinite-index subgroups removed.
If $T_0$ is infinite it is dense in the circle. In a
translation-invariant finitely additive probability on the abelian
group $T_0$, every infinite-index coset has measure zero by disjoint
translates. Every nonempty open arc intersected with $T_0$ has positive
measure, since finitely many translates cover $T_0$. Each such
surviving infinite piece is therefore dense. But the only possible
accumulation point of $D$ is $1$, since $w_\rho\to1$ as the zero
height grows. Thus $D$ is finite. The same Hadamard and shift argument
now gives $\zeta_K=\zeta_L$.
